%#############################################################################
%  模板文件 —— 新建章节时请「复制本文件」再改名，不要直接在 \input 里写正文
%  本文件本身不被主文件调用，可以随意修改，不会影响编译结果
%#############################################################################
%
%  【一】放在哪里
%    chapters\<两位编号>-<章名>.tex
%    例如：chapters\06-平面向量.tex
%          chapters\07-复数.tex
%    编号决定排序，务必补零（01、02……10），不要写成 1、2、10。
%    章名建议与 \chapter{} 里的标题同名，便于查找。
%
%  【二】主文件里要加哪两行（在 讲义.tex 的 \mainmatter 之后）
%    \chapter{平面向量}
%    \input{chapters/06-平面向量}
%    注意：路径用正斜杠 / ，不要用反斜杠 \ ；扩展名 .tex 可写可不写。
%
%  【三】本文件里绝对不能出现的东西
%    \documentclass{...}      导言区已在主文件里，重复写会报错
%    \begin{document} ... \end{document}
%    \chapter{...}            章标题由主文件负责，这里只从 \section 开始
%    \usepackage{...}         宏包统一在主文件导言区添加
%    \setCJKmainfont 等字体命令  同上
%
%  【四】允许直接使用的东西
%    字体、颜色、定理环境、图片搜索路径（figure/、image/）都由主文件提供，
%    这里照常用即可。图片写成 \includegraphics{figure/xxx.png} 或直接写文件名。
%
%#############################################################################


% ======== 以下为示范内容，复制后请整体替换成你自己的正文 ========

\section{平面向量的概念}
\subsection{向量的定义}

\begin{definition}[向量]
既有大小又有方向的量叫做向量。
\end{definition}

% 可选的方括号里写「标题」，会显示为“定义 1.1（向量）”
\begin{definition}[零向量与单位向量]
长度为零的向量叫做零向量，记作 $\vec{0}$；长度等于 $1$ 的向量叫做单位向量。
\end{definition}

\begin{example}
下列各量中是向量的有（\quad）
\begin{enumerate}
    \item 质量
    \item 速度
    \item 温度
    \item 位移
\end{enumerate}
\end{example}

\subsection{向量的表示}

\begin{definition}[向量的几何表示]
用有向线段表示向量，有向线段的长度表示向量的大小，箭头所指的方向表示向量的方向。
\end{definition}

% 定理类环境都可带 [标题]，编号由主文件统一自动递增
\begin{theorem}[向量加法的三角形法则]
已知向量 $\vec{a}$、$\vec{b}$，在平面内任取一点 $A$，作 $\overrightarrow{AB}=\vec{a}$，$\overrightarrow{BC}=\vec{b}$，则 $\overrightarrow{AC}=\vec{a}+\vec{b}$。
\end{theorem}

\begin{lemma}[向量加法的交换律]
对任意向量 $\vec{a}$、$\vec{b}$，都有 $\vec{a}+\vec{b}=\vec{b}+\vec{a}$。
\end{lemma}

\begin{corollary}[向量不等式]
对任意向量 $\vec{a}$、$\vec{b}$，都有 $\bigl||\vec{a}|-|\vec{b}|\bigr| \leq |\vec{a}+\vec{b}| \leq |\vec{a}|+|\vec{b}|$。
\end{corollary}

% 需要插入图片时用 figure 环境；figure/ 和 image/ 已在主文件里设为自动搜索路径
% \begin{figure}[htbp]
%     \centering
%     \includegraphics[width=0.6\textwidth]{figure/示例图.png}
%     \caption{图的说明文字}
%     \label{fig:example}
% \end{figure}

\begin{note}
笔记环境没有编号，适合放易错提醒、补充说明。
\end{note}

\begin{remark}
注记环境用于补充评述，同样没有编号。
\end{remark}

% 带解答的练习题：exercise 自动编号，solution 不编号
\begin{exercise}
已知 $|\vec{a}|=3$，$|\vec{b}|=4$，$\vec{a} \perp \vec{b}$，求 $|\vec{a}+\vec{b}|$。
\end{exercise}
\begin{solution}
因为 $\vec{a} \perp \vec{b}$，所以 $\vec{a}\cdot\vec{b}=0$，故
\[
|\vec{a}+\vec{b}|=\sqrt{|\vec{a}|^2+2\vec{a}\cdot\vec{b}+|\vec{b}|^2}=\sqrt{9+0+16}=5.
\]
\end{solution}

% ======== 示范内容结束 ========
